\section{The cost of limiting independent quantum blocks}
\label{sec:blocklearning78}

The training law changes when the number of calls in one independent
quantum experiment is prescribed.  The relevant restriction concerns
access to the next queried input.  Quantum systems retained from completed
experiments may still be stored and measured together.  This distinction
allows a direct comparison with tomography based on independently
prepared probe--reference pairs.

\begin{definition}[Independent blocks with retained outputs]
\label{def:freshblocks78}
A training protocol has independent blocks of length at most $b$ if it
has the following form.  At a block boundary it may apply an arbitrary
common quantum instrument to all its retained systems.  The resulting
classical history $h$, including its classical random choices and the
number of calls already charged, determines whether to stop, or else a
length $1\le t(h)\le b$ and a fresh $t(h)$-use tester.  Conditional on
$h$, this tester starts in a parameter-independent state tensor-separated
from all retained systems.  It may use arbitrary references, entangled
inputs, and adaptive control within its $t(h)$ calls.

From its preparation until its last call, the entire fresh tester,
including its reference and unused inputs, is isolated from the retained
systems.  Those systems are available again at the next block boundary,
when the new outputs may be joined to them and processed collectively.
Thus classical feedback may select the next tester, while no retained
quantum register can be coupled into a block that still has calls to
make.  A block length is declared before that block and charged in full;
an internal early exit is padded within its declared length.  The sum of
declared lengths is at most $M$ on every record.  Stopping at block
boundaries may depend on the full classical history.  The final estimate
may use an arbitrary joint measurement of all retained systems.
Classical record spaces are standard Borel and quantum registers
are separable.
\end{definition}

The retained quantum systems belong to the tester; the unknown
measurement still consumes its input and returns only its ordered binary
classical label.  The definition permits an unbounded amount of retained
quantum storage and classical computation.  It also permits all probes to be prepared
in advance when their states factor into the prescribed independent
blocks: their preparation can be deferred without changing the
experiment.  In particular, independent one-call Choi preparations
followed by a collective measurement belong to the case $b=1$.
The distinction between classical input control and quantum memory on
the output side has an antecedent in the channel-discrimination model
of \cite[Section~II.2]{SHW68}.  The familiar quadratic information gain
within a block and additive gain across independent blocks also occur
in the metrological bounds of \cite{Hyllus78,Toth78}.  We give a
finite-budget fidelity argument because a local information bound alone
does not supply the confidence dependence needed here.

Write $M_b^\star(d,N,\delta,\eta)$ for the least every-record call
budget of a protocol in Definition~\ref{def:freshblocks78} whose legal
estimate satisfies
\begin{equation}\label{eq:blockrisk78}
 \sup_{E\in\mathfrak E_d}
 \mathbb P_E\{d_N(\mathcal M_E,\mathcal M_{\widehat E})>\delta\}
 \le\eta.
\end{equation}
The future distance $d_N$ retains its unrestricted adaptive definition
and its unhalved trace-norm normalization.  The restriction on $b$
applies only to training.

\begin{theorem}[Sharp tradeoff for independent training blocks]
\label{thm:blocklearning78}
There are absolute constants $c,C,\delta_*>0$ such that, for every
$d\ge2$, $N\ge1$, $1\le b\le N$, $0<\delta\le\delta_*$, and
$0<\eta\le1/8$,
\begin{equation}\label{eq:blockcalls78}
 c\frac{N^2}{b\delta^2}\log\frac1\eta
 \le M_b^\star(d,N,\delta,\eta)
 \le C d^4\frac{N^2}{b\delta^2}\log\frac d\eta.
\end{equation}
The upper bound is attained by one procedure common to the entire
closed effect body.  The lower bound allows classical feedback,
arbitrary quantum storage and collective processing at block
boundaries, and every-record bounded stopping.  Consequently, for
each fixed $d\ge2$,
\begin{equation}\label{eq:blockminimax78}
 M_b^\star(d,N,\delta,\eta)
 =\Theta_d\!\left(\frac{N^2}{b\delta^2}\log\frac1\eta\right).
\end{equation}
In dimension $d=1$ the corresponding order is
$\Theta(N\delta^{-2}\log(1/\eta))$, independently of $b$.
\end{theorem}

The endpoints $b=1$ and $b=N$ differ by a factor of order $N$.
The claim is uniform in the block cap, horizon, accuracy, and confidence;
the upper factor $d^4$ is not asserted to be sharp in growing dimension.

\begin{lemma}[Changing the future horizon]
\label{lem:blockhybrid78}
For all effects $E,F$ and integers $N,b\ge1$,
\begin{equation}\label{eq:blockhybrid78}
 d_N(\mathcal M_E,\mathcal M_F)
 \le\left\lceil\frac Nb\right\rceil
           d_b(\mathcal M_E,\mathcal M_F).
\end{equation}
\end{lemma}
\begin{proof}
Fix an $N$-use tester and partition its calls into consecutive groups
of at most $b$.  Replace the device in one group at a time.  For each
replacement the common prefix prepares a reference-assisted input
to a tester with at most $b$ calls, including all memory carried from
the prefix.  The common suffix is a channel on that tester's output
and cannot increase trace distance.  Each replacement therefore
costs at most $d_b$.  Sum these costs and take the supremum over the
original tester.  Public stopping is included by padding to $N$ calls
with discarded dummy experiments.
\end{proof}

\begin{proof}[Upper bound in Theorem~\ref{thm:blocklearning78}]
Put $q=\lceil N/b\rceil$.  Apply
Theorem~\ref{thm:matrixlearning77} with future horizon $b$, accuracy
$\delta/q$, and failure allowance $\eta$.  Its construction uses fresh
independent blocks of length at most $b$, with classical records between
them, so it satisfies Definition~\ref{def:freshblocks78}.
Lemma~\ref{lem:blockhybrid78} converts its success event into
\eqref{eq:blockrisk78}.  On every record its call count is at most
\[
 C d^4 bq^2\delta^{-2}\log(d/\eta)
 \le 4C d^4\frac{N^2}{b\delta^2}\log(d/\eta),
\]
because $1\le b\le N$.  Choose $\delta_*$ no larger than the
accuracy constant in Theorem~\ref{thm:matrixlearning77}.
\end{proof}

For positive, possibly subnormalized operators define root fidelity by
\[
 \mathfrak f(A,B)=\|\sqrt A\sqrt B\|_1.
\]
We use its tensor-product identity, its additivity on direct sums, and
its monotonicity under a common channel.  In particular, for a common
instrument $(\mathcal J_z)_z$,
\begin{equation}\label{eq:blockinstrument78}
 \sum_z\mathfrak f(\mathcal J_z(A),\mathcal J_z(B))
 \ge\mathfrak f(A,B).
\end{equation}
These facts hold for subnormalized operators by homogeneity.  For
states, the Bures angle $\arccos\mathfrak f$ is a contractive metric;
see \cite{Watrous65,YF76}.

\begin{lemma}[Fidelity with adaptive choices of fresh blocks]
\label{lem:blockfidelity78}
Fix a two-dimensional subspace and, for $\phi\in\mathbb R$, set
\begin{equation}\label{eq:blockprojective78}
 E_\phi=\frac12(I_2+\cos\phi\,\sigma_1+
                         \sin\phi\,\sigma_2)
                   \mathbin\oplus\frac12 I_{d-2},
\end{equation}
with the second summand absent when $d=2$.  If $b|\theta|\le1$,
the final states $\Gamma_0,\Gamma_\theta$ of any protocol in
Definition~\ref{def:freshblocks78}, applied to $E_0,E_\theta$ with
every-record budget $M$, satisfy
\begin{equation}\label{eq:blockfidelity78}
 \mathfrak f(\Gamma_0,\Gamma_\theta)
 \ge\exp(-Mb\theta^2/4).
\end{equation}
Furthermore, if $N|\theta|\le\pi$, then
\begin{equation}\label{eq:blockprojectivedistance78}
 d_N^{\rm na}(\mathcal M_{E_0},\mathcal M_{E_\theta})
 =d_N(\mathcal M_{E_0},\mathcal M_{E_\theta})
 =2\sin(N|\theta|/2).
\end{equation}
\end{lemma}
\begin{proof}
For qubits, \eqref{eq:blockprojectivedistance78} follows from
\cite[Corollary~1 and Theorem~2]{PPKKO72}.  We give the calculation
to specify the normalization and the fidelity estimate needed below.

Let $U_\theta=e^{-\mathrm i\theta\sigma_3/2}\oplus I_{d-2}$.
Then $E_\theta=U_\theta E_0U_\theta^*$ and
$\mathcal M_{E_\theta}=\mathcal M_{E_0}\circ
\operatorname{Ad}_{U_\theta^*}$.  If $|\theta|\le\pi$, every
purified input $\psi$, including an arbitrary reference, satisfies
\[
 \big|\langle\psi,(U_\theta^*\otimes I)\psi\rangle\big|
 \ge \operatorname{Re}\langle\psi,
                    (U_\theta^*\otimes I)\psi\rangle
 \ge\cos(|\theta|/2).
\]
The last inequality follows from the spectrum of
$(U_\theta+U_\theta^*)/2$.  Applying the common measurement and
discarding a purification can only increase fidelity.  Thus one
call changes the output Bures angle by at most $|\theta|/2$,
uniformly over its input and reference.

For a common adaptive tester, apply the triangle inequality before
each query: the angle between the two current states is contracted
by the same channel, and changing the channel on the second state
adds at most $|\theta|/2$.  The interposed controls are common
channels and also contract the angle.  Induction from the common
initial state proves that every $t$-use tester has output states
$\rho_0,\rho_\theta$ satisfying
\begin{equation}\label{eq:blocklocalfidelity78}
 \mathfrak f(\rho_0,\rho_\theta)
 \ge\cos(t|\theta|/2)
 \quad\hbox{when }t|\theta|\le\pi.
\end{equation}
This argument permits arbitrary adaptive quantum memory inside the
tester.  Its output includes all internal outcomes before any
postselection.  For $t\le b$ and $b|\theta|\le1$,
$\cos x\ge e^{-x^2}$ for $0\le x\le1/2$ gives
\begin{equation}\label{eq:blocklocalexponential78}
 \mathfrak f(\rho_0,\rho_\theta)
 \ge e^{-t^2\theta^2/4}\ge e^{-at},
 \qquad a=b\theta^2/4.
\end{equation}

We next retain the whole classical history, including all branch
lengths and stopping decisions.  At a boundary write the two
classical--quantum states as direct sums of subnormalized retained
states $A_h,B_h$.  Their traces include the respective probabilities
of $h$.  Let $m(h)$ be the charged calls in that history, and consider
\begin{equation}\label{eq:blockpotential78}
 \Psi=\sum_h e^{-a(M-m(h))}\mathfrak f(A_h,B_h).
\end{equation}
Initially $\Psi=e^{-aM}$, since the initial resources are the same.
A common boundary instrument preserves $m(h)$ along each of its
children, so \eqref{eq:blockinstrument78} shows that it cannot
decrease $\Psi$.  The history is kept until the end; in particular,
records with different charged call counts are never merged in this
calculation.

At a continuing history $h$, the prescribed fresh tester has the
same specification under both hypotheses.  Conditional on $h$,
its output states $\rho_0^h,\rho_\theta^h$ are normalized and are
tensor-separated from the old retained state.  Immediately after
this block the contribution of this history is therefore
\[
 \begin{split}
 &e^{-a(M-m(h)-t(h))}
   \mathfrak f(A_h\otimes\rho_0^h,
               B_h\otimes\rho_\theta^h)\\
 &\qquad\ge
 e^{-a(M-m(h)-t(h))}e^{-at(h)}\mathfrak f(A_h,B_h)
 =e^{-a(M-m(h))}\mathfrak f(A_h,B_h).
 \end{split}
\]
The subsequent joint instrument is covered by
\eqref{eq:blockinstrument78}.  A stopped history is carried
unchanged.  Since each continuing block charges at least one call,
all histories are terminal after at most $M$ such generations.
At that time all weights in \eqref{eq:blockpotential78} are at
most one.  Additivity of root fidelity on the terminal histories
therefore gives
\[
 \mathfrak f\!\left(\bigoplus_h A_h,\bigoplus_h B_h\right)
 =\sum_h\mathfrak f(A_h,B_h)\ge\Psi\ge e^{-aM}.
\]
Any final collective processing or discarding of history again
increases fidelity.  This proves \eqref{eq:blockfidelity78} without
asserting independence of the unconditional records or imposing a
common support condition on their laws.

For general measurable classical outcomes, choose a measure dominating
the two history laws and interpret $A_h,B_h$ as their operator-valued
densities.  Direct-sum additivity becomes the corresponding integral
identity, and the same proof uses integrals in
\eqref{eq:blockinstrument78} and \eqref{eq:blockpotential78}.
These identities follow equivalently by finite coarse-graining of
the classical record and refinement.  The charged call count has
only the values $0,\ldots,M$ and is retained in every such
coarse-graining.

Finally, use the first two coordinates for a parallel $N$-input state
\[
 \frac{|0\rangle^{\otimes N}
             +e^{\mathrm i\alpha}|1\rangle^{\otimes N}}{\sqrt2},
 \qquad \alpha=N\theta/2+\pi/2.
\]
The parity of the ordered output signs has mean
$\cos(N\phi-\alpha)$ under $E_\phi$, hence means
$-\sin(N\theta/2)$ and $\sin(N\theta/2)$ under the two
hypotheses.  Its unhalved distance is $2|\sin(N\theta/2)|$.
For $N|\theta|\le\pi$, apply
\eqref{eq:blocklocalfidelity78} with $t=N$ to any future adaptive
tester and use
$\|\rho-\sigma\|_1\le2\sqrt{1-\mathfrak f(\rho,\sigma)^2}$.
The resulting upper bound is $2\sin(N|\theta|/2)$, equal to the
parallel witness.  This proves \eqref{eq:blockprojectivedistance78}.
\end{proof}

\begin{proof}[Lower bound in Theorem~\ref{thm:blocklearning78}]
Reduce $\delta_*$ further so that $\delta_*\le1/4$, and set
$\theta=4\delta/N$ in \eqref{eq:blockprojective78}.
Then $b|\theta|\le1$, while
\eqref{eq:blockprojectivedistance78} gives
\[
 d_N(\mathcal M_{E_0},\mathcal M_{E_\theta})
 =2\sin(2\delta)>2\delta.
\]
The two radius-$\delta$ success balls are disjoint.  Thus any
estimate satisfying \eqref{eq:blockrisk78} determines a test with
both error probabilities $\alpha,\beta$ at most $\eta$: choose
the first hypothesis precisely when the estimate lies in its
radius-$\delta$ ball.  The classical root fidelity of that test is
\[
 \sqrt{(1-\alpha)\beta}+\sqrt{\alpha(1-\beta)}
 \le2\sqrt\eta.
\]
Fidelity is nondecreasing under the map from the training output to
this test.  Lemma~\ref{lem:blockfidelity78} consequently implies
\[
 e^{-Mb\theta^2/4}\le2\sqrt\eta,
 \qquad
 M\ge\frac{N^2}{8b\delta^2}\log\frac1{4\eta}
 \ge\frac{N^2}{24b\delta^2}\log\frac1\eta,
\]
where the last inequality uses $\eta\le1/8$.
This proves the asserted lower bound in every $d\ge2$.

For $d=1$ every effect is scalar.  Bernoulli endpoint estimation
in Lemma~\ref{lem:learnendpoint76} gives the upper order with
independent one-call samples.  The scalar converse in
Theorem~\ref{thm:matrixlearning77} gives the matching lower order
even with unrestricted adaptive quantum training.  Thus increasing
$b$ has no effect in dimension one.
\end{proof}

\begin{corollary}[The limitation of independent one-call acquisition]
\label{cor:onecallseparation78}
Fix $d\ge2$.  Every learner that acquires its data using fresh
one-call probes and fresh references, with classical feedback and
arbitrary subsequent collective processing as in
Definition~\ref{def:freshblocks78}, requires
$\Omega_d(N^2\delta^{-2}\log(1/\eta))$ calls for
\eqref{eq:blockrisk78}.  The complete closed effect body admits
the same risk with $O_d(N\delta^{-2}\log(1/\eta))$ calls when
independent blocks of length up to $N$ are permitted.
\end{corollary}
\begin{proof}
Use $b=1$ and $b=N$ in Theorem~\ref{thm:blocklearning78}.
\end{proof}

This conclusion applies to the acquisition interfaces of
\cite{MB67,ZRK76}.  The channel learner in
\cite[Section~III.2 and Algorithm~I.2]{MB67} prepares independent
Choi copies before their collective processing.  The measurement
tomography procedure in \cite[Section~3]{ZRK76} queries separately
prepared known inputs.  Both are contained in $b=1$ above.  Any
estimator retaining these interfaces and required to satisfy the
present worst-case future-$N$ loss therefore incurs the quadratic
horizon order, including when its final estimator is collective.
The block lower bound settles this interface comparison directly;
it does not rely on substituting $\delta/N$ into a sufficient
one-use accuracy guarantee.

\begin{corollary}[Learning with bounded blocks into an optimal-order word]
\label{cor:blocklearnedcode78}
For fixed $d\ge2$, $1\le b\le N$, rational
$0<\delta\le\min\{1,\delta_*,\delta_d\}$, and
$0<\eta\le1/8$, there is a common learner in
Definition~\ref{def:freshblocks78} that returns a fixed-length word
whose deterministic public decoder produces a legal effect $C$ with
\[
 \sup_{E\in\mathfrak E_d}
 \mathbb P_E\{d_N(\mathcal M_E,\mathcal M_C)>\delta\}\le\eta.
\]
Its every-record training budget is at most
$Cd^4N^2(b\delta^2)^{-1}\log(d/\eta)$ and its payload has length
\begin{equation}\label{eq:blockcodebits78}
 \frac{d^2}{2}\log_2N+
 \lfloor d/2\rfloor\log_2\log(N+2)+
 d^2\log_2(1/\delta)+O_d(1).
\end{equation}
For each fixed $d$, both the training order in this block model and
the payload order among fixed public decoders into legal effects are
optimal.  Encoding uses no additional device calls.
\end{corollary}
\begin{proof}
Run the learner of Theorem~\ref{thm:blocklearning78} at accuracy
$\delta/2$.  The upper construction inherits the legal rational
output of Theorem~\ref{thm:matrixlearning77}.  Apply the exact
codec of Theorem~\ref{thm:matrixcodec77} at accuracy $\delta/2$.
On the learning success event its decoded effect obeys
\[
 d_N(\mathcal M_E,\mathcal M_C)
 \le\delta/2+11\delta/32=27\delta/32<\delta.
\]
The codec is a deterministic computation on the recorded rational
estimate.  Theorem~\ref{thm:matrixcodec77} gives
\eqref{eq:blockcodebits78}.  The training converse is
Theorem~\ref{thm:blocklearning78}.  For the payload converse,
positive success probability for every $E$ forces the decoded
family to cover $\mathfrak E_d$ by radius-$\delta$ balls; apply
Theorem~\ref{thm:matrixcover76}.  As before, reconstruction time,
workspace, and state-preparation complexity are separate from the
two resource counts.
\end{proof}
